%! The Normal Distribution and z-Scores — Unit 1, Lesson 1.7 — Warm-Up
% ONE PAGE, blank and key. Three items. 12pt, like every student component.
% CONVERTED FROM 10pt: going to 12pt costs roughly 20% of the vertical budget, so the structural
% spacing was trimmed (itemsep 0.19in -> 0.13in, the \vspace values, the \blank widths, the
% table's \arraystretch) and item 2's first two sub-parts were merged into one. Budget is
% \pageheader + three short items — if anything is added here, something else has to go.
% Mean and standard deviation were formalized in Lesson 1.5, the five-number summary and IQR in
% Lesson 1.6 — prior knowledge, used freely.
% Item 2 seeds the z-score arithmetic of notes row 2, done once with numbers that divide evenly;
% item 3 seeds the halve-the-tail move that notes problems 23-24 rest on. Item 3's percentages
% are 90/10/5 rather than the real 95/5/2.5 so that the empirical rule's own numbers arrive out
% of the students' band counts in notes problem 13.
\documentclass[12pt]{article}
\usepackage{apstats-article}
\usepackage{apstats-boxes}

\begin{document}

\pageheader{Unit 1, Lesson 1.7}{Warm-Up}

\vspace{0.12in}

\begin{enumerate}[leftmargin=1.7em, itemsep=0.13in]

  \item \textbf{From last lesson.} Here are five-number summaries for the daily commute times,
        in minutes, of two groups of workers.

        \vspace{3pt}
        \renewcommand{\arraystretch}{1.25}
        \begin{tabularx}{\linewidth}{@{} l Y Y Y Y Y @{}}
        \rowcolor{sky}
        \textbf{Group} & \textbf{Min} & \textbf{Q1} & \textbf{Med} & \textbf{Q3} & \textbf{Max} \\
        \hline
        \textbf{Downtown} & 8 & 14 & 19 & 26 & 41 \\
        \textbf{Suburban} & 15 & 22 & 25 & 30 & 38 \\
        \end{tabularx}
        \vspace{4pt}

        IQR for Downtown: \blank{2.0cm} \qquad IQR for Suburban: \blank{2.0cm}

        \vspace{3pt}
        Whose \emph{middle half} of commute times is more spread out? \blank{3.7cm}

  \item A set of quiz scores has a mean of \textbf{20} points and a standard deviation of
        \textbf{3} points.

        \vspace{3pt}
        \hspace{1.0em}\textbf{(a)} A score of 26 is how many points above the mean, and how many
        3-point steps is that? \blank{3.6cm}
        \vspace{3pt}

        \hspace{1.0em}\textbf{(b)} A score of 17 is how many 3-point steps from the mean, and on
        which side? \blank{4.3cm}

  \item In a set of \textbf{400} measurements, \textbf{360} of them fall inside a certain range
        of values.

        \vspace{3pt}
        \hspace{1.0em}\textbf{(a)} What percent of the measurements is that? \blank{2.0cm}
        \vspace{3pt}

        \hspace{1.0em}\textbf{(b)} What percent fall \emph{outside} that range? \blank{2.0cm}
        \vspace{3pt}

        \hspace{1.0em}\textbf{(c)} The measurements pile up evenly on both sides of the middle,
        so the leftovers split evenly between the low end and the high end. What percent fall
        past the \emph{high} end? \blank{2.0cm}

\end{enumerate}

\end{document}
